\documentclass[aspectratio=169]{beamer}
\usetheme{metropolis}
\metroset{sectionpage=none, numbering=counter, progressbar=frametitle}
\usepackage{amsmath, amssymb}
\usepackage[T1]{fontenc}

\title{Domain-Partition Algorithm for Turbomachinery Blades}
\subtitle{Slide 11: Periodic Boundaries --- 3$\times$ Tiled Passages}
\date{\today}
\author{}
\institute{}

\begin{document}

\begin{frame}{Periodic Boundaries --- 3$\times$ Tiled Passages}

\begin{itemize}
    \item \textbf{3 passages tiled:} shifts $-\text{pitch}$, $0$, $+\text{pitch}$
    \item \textbf{Interior seams:} grid points of neighboring copies coincide (CFD-conforming)
    \item \textbf{Visual check:} master-slave conformity confirmed by overlay
\end{itemize}

\vspace{0.5em}
\small
\textit{Code:} \texttt{tmesh\_partition.py:1271--1312} (\texttt{plot\_tiled})
\hfill
\textit{Lit:} DLR Sauer/Morsbach 2023, sec 2.7 master-slave

\note{
Speaker notes:
The plot_tiled function in tmesh_partition.py renders three hub passages side by side, shifted by minus pitch, zero, and plus pitch. This is the visual periodicity check the user asked for. At the interior seams between the center passage and its left and right copies, the grid points must coincide exactly. This coincidence is what makes the mesh CFD-conforming for periodic boundary conditions. The green seam curves overlaid on the plot confirm that the master and slave boundaries align. The DLR Sauer and Morsbach 2023 paper, section 2.7, describes the master-slave pairing formalism that underlies this check.
}

\end{frame}

\end{document}
